\section{Related Work}
\label{sec:related_work}

\subsection{Image Intent Recognition}
Intent recognition seeks to decode the underlying purposes of humans. 
Early studies mainly focus on specific domains. \revised{Joo et al.~\cite{joo_visual_2014} infer communicative intent from persuasive images through facial and contextual cues, and subsequently relate facial-trait judgments to election outcomes~\cite{joo_automated_2015}; Huang and Kovashka~\cite{huang_inferring_2016} further infer visual persuasion by jointly modeling body language, scene setting, and deep features.}
More recently, the release of the Intentonomy dataset \cite{jia_intentonomy_2021} has facilitated the study of general-purpose visual intention understanding in unconstrained social media scenarios.

Existing methods for visual intention recognition seek to solve the problem from different perspectives. 
Methods \textit{e.g.} PIP-Net \cite{wang_prototype-based_2023} concentrate on feature and representation learning, which introduce prototype-based representations to capture diverse visual patterns and improve interpretability. 
Considering the hierarchical nature of intent labels, approaches like CPAD \cite{ye_cross-modality_2023} and HLEG \cite{shi_learnable_2023} exploit label hierarchies to improve multi-level semantic alignment and mitigate ambiguity.  
LabCR \cite{shi_label-aware_2024} addresses label noise and shifting through calibration strategies, while recent works including IntCLIP \cite{leonardis_synergy_2025} and IntentMLM \cite{wang_unlocking_2026} leverage vision-language models and textual priors to align visual features with semantic intent descriptions.

Despite these advances, most existing methods address ambiguity from a single perspective, such as feature representation, label structure, or external knowledge. 
They implicitly assume that ambiguity should be handled uniformly within a model. 
However, in visual intention recognition, ambiguity arises from different sources and may affect different stages of the recognition process in distinct ways.

\subsection{Learning under Ambiguous or Uncertain Supervision}

Learning under ambiguous or noisy supervision has been extensively studied in multi-label recognition and subjective understanding tasks, where label reliability is often compromised by annotation noise, semantic overlap, or annotator disagreement. 
Existing approaches typically address this issue in three main directions.

First, noise-robust learning methods mitigate unreliable supervision through robust objectives and example reweighting: \revised{Focal Loss~\cite{lin_focal_2020} down-weights easy examples to focus training on hard ones, online hard-example mining~\cite{shrivastava_training_2016} explicitly selects difficult samples during detector training, collaborative refining~\cite{ye_collaborative_2022} suppresses label noise in person re-identification, and the asymmetric loss~\cite{ridnik_asymmetric_2021} re-balances positive and negative gradients for multi-label classification.}
These methods reduce the adverse impact of corrupted or hard labels, but they usually treat uncertainty as a generic optimization issue rather than modeling its semantic origin. More recently, several studies have further argued that disagreement itself can be informative rather than purely harmful. Instead of collapsing subjective annotations into a single majority-vote label, these works preserve human uncertainty or annotator-level diversity as part of the learning signal: \revised{Peterson et al.~\cite{peterson_human_2019} show that training on human label distributions makes classification more robust, Davani et al.~\cite{mostafazadeh_davani_dealing_2022} model individual annotators rather than the majority vote, DisCo~\cite{weerasooriya_disagreement_2023} jointly models item- and annotator-level label distributions, and CrowdOpinion~\cite{weerasooriya_subjective_2023} applies population-level learning to recover meaningful minority opinions.} This perspective is especially relevant to visual intention recognition, where multiple interpretations are often simultaneously plausible.

Second, hierarchical or structure-aware methods improve learning under ambiguous labels by exploiting label relations or taxonomy constraints \cite{zhang_use_2022, shi_learnable_2023}. 
Such methods are especially relevant to intention recognition, where semantically nearby labels often exhibit overlapping decision boundaries. 
However, they mainly focus on label dependency modeling, instead of explicitly using annotator disagreement to regulate supervision strength.

Third, knowledge distillation and vision-language supervision have been adopted to inject additional semantic guidance from teacher models or pretrained language priors \cite{touvron_training_2021, radford_learning_2021}. 
Knowledge distillation transfers structured knowledge from stronger teacher models, while vision-language models leverage textual descriptions as semantic anchors to enhance visual representations. 
However, these approaches typically assume that the provided semantic guidance is uniformly reliable across samples, which may not hold in subjective tasks with inherent ambiguity.

In visual intention understanding, ambiguity is intrinsic: the same image may support multiple plausible interpretations, and category boundaries are often fuzzy \cite{jia_intentonomy_2021, shi_learnable_2023}. 
Prior work has recognized the importance of hierarchical relations and label ambiguity in this task, but existing methods still largely rely on fixed supervision or globally applied semantic guidance. 
By contrast, our approach explicitly leverages annotator disagreement as an uncertainty signal, and adaptively regulates semantic guidance according to supervision ambiguity. 
This results in a more robust learning process under subjective and ambiguous supervision.